% Part: methods
% Chapter: proofs
% Section: example-1

\documentclass[../../../include/open-logic-section]{subfiles}

\begin{document}

\olfileid{mth}{prf}{ex1}

\olsection{એક ઉદાહરણ}

આપણું પહેલું ઉદાહરણ ગણોના યોગગણ અને છેદગણ અંગેની નીચેની
સરળ હકીકત છે. તે વ્યાખ્યાઓ ઉકેલવી, સંયોજન અને સાર્વત્રિક દાવા
સાબિત કરવા તથા કિસ્સાવાર સાબિતી કરવી બતાવશે.

\begin{prop}
કોઈ પણ ગણો $A$, $B$ અને $C$ માટે $A \cup (B \cap C) = (A \cup B)
\cap (A \cup C)$.
\end{prop}

ચાલો તે સાબિત કરીએ!{}

\begin{proof}
આપણે બતાવવું છે કે કોઈ પણ ગણો $A$, $B$ અને $C$ માટે $A \cup (B \cap
C) = (A \cup B) \cap (A \cup C)$.
\begin{quote}
પહેલાં પ્રતિજ્ઞાપનમાંના ``$=$''ની વ્યાખ્યા ઉકેલીએ. યાદ રાખો કે
ગણો એકસમાન છે તે સાબિત કરવા તેમના !!{element}s સરખા છે તે
બતાવવું પડે. એટલે $A \cup
(B \cap C)$નાં બધાં !!{element}s, $(A \cup B) \cap (A \cup C)$નાં
પણ !!{element}s છે અને ઊલટું પણ છે. ``ઊલટું પણ''નો અર્થ
એ છે કે $(A
\cup B) \cap (A \cup C)$નું દરેક !!{element}, $A \cup (B \cap
C)$નું !!a{element} પણ હોવું જોઈએ. તેથી વ્યાખ્યા ઉકેલતાં
આપણે સંયોજન સાબિત કરવાનું છે તે સમજાય છે. તેને નોંધીએ:
\end{quote}
વ્યાખ્યા મુજબ, $A \cup (B \cap C) = (A \cup B) \cap (A \cup C)$
ત્યારે અને ત્યારે જ જ્યારે $A \cup (B \cap C)$નું દરેક
!!{element}, $(A
\cup B) \cap (A \cup C)$નું !!a{element} પણ છે અને
$(A \cup B) \cap (A
\cup C)$નું દરેક !!{element}, $A \cup (B \cap C)$નું
!!a{element} પણ છે.
\begin{quote}
આ સંયોજન હોવાથી તેના બંને ભાગો અલગથી સાબિત કરવા પડે.
પહેલાથી શરૂ કરીએ: $A \cup (B \cap C)$નું દરેક
!!{element}, $(A
\cup B) \cap (A \cup C)$નું !!a{element} પણ છે તે સાબિત કરીએ.

આ સાર્વત્રિક દાવો છે; તેથી $A \cup (B \cap C)$નું મનસ્વી !!{element}
લઈએ અને બતાવીએ કે તે $(A \cup B) \cap (A \cup C)$નું
!!a{element} પણ હોવું જોઈએ. આ મનસ્વી !!{element}ને કોઈ
ચલથી ઓળખાવીએ, ધારો કે~$z$. હવે સાબિતી આગળ વધે છે:
\end{quote}
પહેલાં આપણે સાબિત કરીએ કે $A \cup (B \cap C)$નું દરેક
!!{element}, $(A \cup B) \cap (A \cup C)$નું !!a{element}
પણ છે. $z \in A \cup (B
\cap C)$ લો. આપણે $z \in (A \cup B) \cap (A \cup C)$ બતાવવું છે.
\begin{quote}
હવે $\cup$ અને~$\cap$ની વ્યાખ્યાઓ ઉકેલવાનો સમય છે. દાખલા તરીકે,
$\cup$ની વ્યાખ્યા $A \cup B = \Setabs{z}{z \in A
  \text{ અથવા } z \in B}$ છે. આ વ્યાખ્યા ``$A \cup (B
\cap C)$''ને લાગુ પાડીએ ત્યારે વ્યાખ્યાના ``$B$''નું સ્થાન
``$B \cap C$'' લે છે; તેથી $A \cup (B \cap C) = \Setabs{z}{z \in A \text{ અથવા }
  z \in B \cap C}$. એટલે $z \in A \cup (B \cap C)$
એવી આપણી પૂર્વધારણાનો અર્થ $z \in \Setabs{z}{z \in A \text{ અથવા } z \in B \cap
  C}$ થાય છે. અને $z \in \Setabs{z}{\dots z\dots}$ ત્યારે અને
ત્યારે જ જ્યારે \dots $z$ \dots; આ કિસ્સામાં $z \in A$
અથવા $z \in B \cap C$.
\end{quote}
$\cup$ની વ્યાખ્યા મુજબ, કાં તો $z \in A$ અથવા $z \in B \cap C$.
\begin{quote}
આ વિયોજન હોવાથી કિસ્સાવાર સાબિતી ઉપયોગી થશે. બે કિસ્સા લઈએ
અને દરેકમાં આપણો લક્ષિત નિષ્કર્ષ (એટલે ``$z \in (A \cup B) \cap (A \cup
C)$'') મળે છે તે બતાવીએ.
\end{quote}
કિસ્સો 1: ધારો કે $z \in A$.
\begin{quote}
આ પૂર્વધારણાઓમાંથી વધુ બહુ કંઈ કાઢવાનું નથી. એટલે નિષ્કર્ષમાં
આપણી પાસે શું છે તે જોઈએ. આપણે $z \in (A \cup B) \cap (A \cup C)$
બતાવવું છે. $\cap$ની વ્યાખ્યા મુજબ, $z \in (A \cup B) \cap (A \cup C)$
બતાવવા તે $(A \cup B)$ અને $(A \cup C)$ બંનેમાં છે તે બતાવવું
પડે. પરંતુ $z
\in A \cup B$ ત્યારે અને ત્યારે જ જ્યારે $z \in A$ અથવા $z \in B$;
અને કિસ્સો~1ની પૂર્વધારણા મુજબ $z \in A$ છે. એ જ
રીતે---$B$ના સ્થાને $C$ મૂકતાં---$z \in A \cup C$ મળે છે.
અહીં વિચાર ઊલટી દિશામાં થયો છે; તેથી સાબિતીમાં જરૂરી દિશામાં
આ તર્ક લખીએ.
\end{quote}
$z \in A$ હોવાથી $z \in A$ અથવા $z \in B$; એટલે~$\cup$ની
વ્યાખ્યા મુજબ $z \in A \cup B$. એ જ રીતે $z \in A \cup C$.
પરંતુ~$\cap$ની વ્યાખ્યા મુજબ તેનો અર્થ $z \in (A \cup B) \cap (A \cup C)$ થાય છે.
\begin{quote}
આથી કિસ્સાવાર સાબિતીનો પહેલો કિસ્સો પૂરો થયો. હવે બીજા
કિસ્સામાં, જ્યાં $z \in B \cap C$ છે, ત્યાં પણ નિષ્કર્ષ મેળવવો છે.
\end{quote}
કિસ્સો 2: ધારો કે $z \in B \cap C$.
\begin{quote}
ફરી આપણે બે ગણોના છેદગણ સાથે કામ કરીએ છીએ. $\cap$ની
વ્યાખ્યા લાગુ પાડીએ:
\end{quote}
$z \in B \cap C$ હોવાથી~$\cap$ની વ્યાખ્યા મુજબ $z$, $B$
અને $C$ બંનેનું !!a{element} હોવું જોઈએ.
\begin{quote}
હવે ફરી નિષ્કર્ષ જોઈએ. $z$ એ $(A \cup B)$ અને $(A \cup C)$
બંનેમાં છે તે બતાવવું છે. ફરી જવાબ તરત મળે છે.
\end{quote}
$z \in B$ હોવાથી $z \in (A \cup B)$. $z \in C$ હોવાથી
$z \in (A
\cup C)$ પણ છે. તેથી $z \in (A \cup B) \cap (A \cup C)$.
\begin{quote}
અહીં ફરી $\cup$ અને $\cap$ની વ્યાખ્યાઓ લાગુ પાડી છે.
પરંતુ આપણે તે પહેલેથી યાદ કરી છે અને $z$ બેમાંથી કોઈ એક
ગણમાં હોય તો તેમના યોગગણમાં પણ છે તે અગાઉ બતાવ્યું છે;
તેથી આ પગલું એટલું સ્પષ્ટ લખવાની જરૂર નથી.

કિસ્સાવાર સાબિતીનો બીજો કિસ્સો પણ પૂરો થયો; હવે પહેલો
નિષ્કર્ષ કહી શકીએ.
\end{quote}
આથી જો $z \in A \cup (B \cap C)$, તો $z \in (A \cup B) \cap (A \cup C)$.
\begin{quote}
હવે વિપરીત દિશા બતાવવાની છે: $(A \cup B) \cap (A \cup C)$નું
દરેક !!{element}, $A \cup (B \cap
C)$નું !!a{element} છે. પહેલાંની જેમ, પહેલા ગણનું મનસ્વી
!!{element} લઈને તે બીજા ગણમાં પણ છે તે બતાવીએ. હવે શું
કરવાના છીએ તે કહીએ.
\end{quote}
હવે ધારો કે $z \in (A \cup B) \cap (A \cup C)$. આપણે
$z \in A \cup (B \cap C)$ બતાવવું છે.
\begin{quote}
હવે $z \in (A \cup B) \cap (A
\cup C)$ એવી પૂર્વધારણાથી કામ કરીએ છીએ. સાબિતીના પહેલા ભાગ
જેવું જ~$z$ અહીં વાપરીએ છીએ, તેથી ગૂંચવણ ન થવી જોઈએ. પહેલો
ભાગ પૂરો કરતાં ત્યાંની બધી પૂર્વધારણાઓ હવે લાગુ નથી પડતી;
તેથી $z$ વિશે નવી પૂર્વધારણાઓ લઈ શકીએ. ગૂંચવણ થાય તો આગળના
ભાગમાં $z$ના સ્થાને બીજો ચલ વાપરો.

$\cap$ની વ્યાખ્યા મુજબ $z$, $A \cup B$ અને $A \cup C$
બંનેમાં છે. $\cup$ની વ્યાખ્યા ઉકેલતાં વધુમાં મળે છે: કાં તો
$z \in A$ અથવા $z \in B$, અને સાથે કાં તો $z \in A$
અથવા $z \in
C$. આ ફરી કિસ્સાવાર સાબિતી જેવું લાગે છે---પણ વચ્ચેનું
``અને'' ગૂંચવે છે. તમને લાગે કે માત્ર ત્રણ શક્યતાઓ છે:
$z$ એ $A$, $B$ કે $C$માંથી કોઈ એકમાં છે. પરંતુ તે ભૂલ
હશે. સાવચેતીથી દરેક વિયોજન અલગથી જોઈએ.
\end{quote}
$\cap$ની વ્યાખ્યા મુજબ $z \in A \cup B$ અને $z \in A \cup C$.
$\cup$ની વ્યાખ્યા મુજબ $z \in A$ અથવા $z \in B$.
કિસ્સા અલગ પાડીએ.
\begin{quote}
હાલમાં પહેલું વિયોજન લઈએ છીએ, તેથી $A \cup C$ ઉકેલીને
મળતું બીજું વિયોજન હજી વાપર્યું નથી. હકીકતમાં તેની હજી જરૂર
પણ નથી. પહેલો કિસ્સો $z \in A$ છે, અને કોઈ ગણનું
!!a{element} તે ગણના બીજા કોઈ પણ ગણ સાથેના યોગગણનું પણ
!!a{element} છે. તેથી કિસ્સો~1 સરળ છે:
\end{quote}
કિસ્સો 1: ધારો કે $z \in A$. તેથી $z \in A \cup (B \cap
C)$.
\begin{quote}
હવે બીજો કિસ્સો, $z \in B$. અહીં બીજો $\cup$ ઉકેલીને
ફરી કિસ્સા અલગ પાડીશું:
\end{quote}
કિસ્સો 2: ધારો કે $z \in B$. $z \in A \cup C$ હોવાથી
કાં તો $z \in
A$ અથવા $z \in C$. કિસ્સા વધુ અલગ પાડીએ:

કિસ્સો 2a: $z \in A$. તો ફરી $z \in A \cup (B \cap C)$.
\begin{quote}
આ થોડું વિચિત્ર હતું. આ કિસ્સામાં~$z \in B$ એવી
પૂર્વધારણાની ખરેખર જરૂર પડી નહીં; તેમાં વાંધો નથી.
\end{quote}
કિસ્સો 2b: $z \in C$. ત્યારે $z \in B$ અને $z \in C$;
તેથી $z \in B \cap C$ અને પરિણામે $z \in A \cup (B \cap C)$.
\begin{quote}
આથી બંને કિસ્સાવાર સાબિતીઓ પૂરી થાય છે અને વિપરીત દિશા
પણ સાબિત થાય છે.
\end{quote}
તેથી જો $z \in (A \cup B) \cap (A \cup C)$ તો $z \in A \cup (B \cap C)$.
\end{proof}

\end{document}
